% GENERATED FILE. Edit canonical lessons, then run npm run generate:books.
% canonical-source-hash: fnv1a64:67560a088e149bf1
% canonical-lessons: BN-C148-ranna-kora, BN-C148-selai-kora, BN-C148-istri-kora, BN-C148-jhat-deoya, BN-C148-seddho-kora

\chapter{To cook, To sew, To iron, To sweep, To boil}
\label{ch:bn-to-cook-to-sew-to-iron-to-sweep-to-boil}
\hlchaptermodality{148}

\begin{chapteropening}
\textbf{By the end of this chapter:} \emph{I can say to cook, to sew, to iron, to sweep and to boil.}

Complete the last lesson of chapter 148: I can say to cook, to sew, to iron, to sweep and to boil.
\end{chapteropening}

\section[\texorpdfstring{\bn{রান্না} \bn{করা} (rānnā kôrā)}{রান্না করা (rānnā kôrā)}]{\bn{রান্না} \bn{করা} (rānnā kôrā) — to cook}

\label{lesson:BN-C148-ranna-kora}



\begin{quote}
Before the new one: say the Bengali for an advert, then the Bengali for a receipt.
\end{quote}

\subsection*{You'll want to know: \bn{রান্না} \bn{করা}}
\textbf{\bn{রান্না} \bn{করা}} — \emph{rānnā kôrā} — \textquotedblleft{}to cook\textquotedblright{}.

\begin{grammarlens}[title={What you've built}]
One more word for this chapter.
\end{grammarlens}

\subsection*{Your turn}
\begin{itemize}
\raggedright
  \item \emph{Say it:} \emph{rānnā kôrā}
  \item \emph{Say it:} \emph{rānnā kôrā}, once more
  \item \emph{Say it:} say \emph{rānnā kôrā}
  \item \emph{Recall:} say the Bengali for an advert, then the Bengali for a receipt, then say \emph{rānnā kôrā} again
\end{itemize}

\begin{tcolorbox}[breakable,colback=teal!4,colframe=teal!35!black,title={Before you move on}]
What does \bn{রান্না} \bn{করা} mean? (\textquotedblleft{}To cook\textquotedblright{}.) Say it once more.
\end{tcolorbox}
\section[\texorpdfstring{\bn{সেলাই} \bn{করা} (selāi kôrā)}{সেলাই করা (selāi kôrā)}]{\bn{সেলাই} \bn{করা} (selāi kôrā) — to sew}

\label{lesson:BN-C148-selai-kora}



\begin{quote}
Before the new one: say the Bengali for a receipt, then the Bengali for to cook.
\end{quote}

\subsection*{You'll want to know: \bn{সেলাই} \bn{করা}}
\textbf{\bn{সেলাই} \bn{করা}} — \emph{selāi kôrā} — \textquotedblleft{}to sew\textquotedblright{}.

\begin{grammarlens}[title={What you've built}]
One more word for this chapter.
\end{grammarlens}

\subsection*{Your turn}
\begin{itemize}
\raggedright
  \item \emph{Say it:} \emph{selāi kôrā}
  \item \emph{Say it:} \emph{selāi kôrā}, once more
  \item \emph{Say it:} say \emph{selāi kôrā}
  \item \emph{Recall:} say the Bengali for a receipt, then the Bengali for to cook, then say \emph{selāi kôrā} again
\end{itemize}

\begin{tcolorbox}[breakable,colback=teal!4,colframe=teal!35!black,title={Before you move on}]
What does \bn{সেলাই} \bn{করা} mean? (\textquotedblleft{}To sew\textquotedblright{}.) Say it once more.
\end{tcolorbox}
\section[\texorpdfstring{\bn{ইস্ত্রি} \bn{করা} (istri kôrā)}{ইস্ত্রি করা (istri kôrā)}]{\bn{ইস্ত্রি} \bn{করা} (istri kôrā) — to iron}

\label{lesson:BN-C148-istri-kora}



\begin{quote}
Before the new one: say the Bengali for to cook, then the Bengali for to sew.
\end{quote}

\subsection*{You'll want to know: \bn{ইস্ত্রি} \bn{করা}}
\textbf{\bn{ইস্ত্রি} \bn{করা}} — \emph{istri kôrā} — \textquotedblleft{}to iron\textquotedblright{}.

\begin{grammarlens}[title={What you've built}]
One more word for this chapter.
\end{grammarlens}

\subsection*{Your turn}
\begin{itemize}
\raggedright
  \item \emph{Say it:} \emph{istri kôrā}
  \item \emph{Say it:} \emph{istri kôrā}, once more
  \item \emph{Say it:} say \emph{istri kôrā}
  \item \emph{Recall:} say the Bengali for to cook, then the Bengali for to sew, then say \emph{istri kôrā} again
\end{itemize}

\begin{tcolorbox}[breakable,colback=teal!4,colframe=teal!35!black,title={Before you move on}]
What does \bn{ইস্ত্রি} \bn{করা} mean? (\textquotedblleft{}To iron\textquotedblright{}.) Say it once more.
\end{tcolorbox}
\section[\texorpdfstring{\bn{ঝাঁট} \bn{দেওয়া} (jhā̃ṭ deoyā)}{ঝাঁট দেওয়া (jhā̃ṭ deoyā)}]{\bn{ঝাঁট} \bn{দেওয়া} (jhā̃ṭ deoyā) — to sweep}

\label{lesson:BN-C148-jhat-deoya}



\begin{quote}
Before the new one: say the Bengali for to sew, then the Bengali for to iron.
\end{quote}

\subsection*{You'll want to know: \bn{ঝাঁট} \bn{দেওয়া}}
\textbf{\bn{ঝাঁট} \bn{দেওয়া}} — \emph{jhā̃ṭ deoyā} — \textquotedblleft{}to sweep\textquotedblright{}.

\begin{grammarlens}[title={What you've built}]
One more word for this chapter.
\end{grammarlens}

\subsection*{Your turn}
\begin{itemize}
\raggedright
  \item \emph{Say it:} \emph{jhā̃ṭ deoyā}
  \item \emph{Say it:} \emph{jhā̃ṭ deoyā}, once more
  \item \emph{Say it:} say \emph{jhā̃ṭ deoyā}
  \item \emph{Recall:} say the Bengali for to sew, then the Bengali for to iron, then say \emph{jhā̃ṭ deoyā} again
\end{itemize}

\begin{tcolorbox}[breakable,colback=teal!4,colframe=teal!35!black,title={Before you move on}]
What does \bn{ঝাঁট} \bn{দেওয়া} mean? (\textquotedblleft{}To sweep\textquotedblright{}.) Say it once more.
\end{tcolorbox}
\section[\texorpdfstring{\bn{সেদ্ধ} \bn{করা} (seddhô kôrā)}{সেদ্ধ করা (seddhô kôrā)}]{\bn{সেদ্ধ} \bn{করা} (seddhô kôrā) — to boil}

\label{lesson:BN-C148-seddho-kora}



\begin{quote}
Before the new one: say the Bengali for to iron, then the Bengali for to sweep.
\end{quote}

\subsection*{You'll want to know: \bn{সেদ্ধ} \bn{করা}}
\textbf{\bn{সেদ্ধ} \bn{করা}} — \emph{seddhô kôrā} — \textquotedblleft{}to boil\textquotedblright{}.

\begin{grammarlens}[title={What you've built}]
One more word for this chapter.
\end{grammarlens}

\subsection*{Your turn}
\begin{itemize}
\raggedright
  \item \emph{Say it:} \emph{seddhô kôrā}
  \item \emph{Say it:} \emph{seddhô kôrā}, once more
  \item \emph{Say it:} say \emph{seddhô kôrā}
  \item \emph{Recall:} say the Bengali for to iron, then the Bengali for to sweep, then say \emph{seddhô kôrā} again
\end{itemize}

\begin{tcolorbox}[breakable,colback=teal!4,colframe=teal!35!black,title={Before you move on}]
What does \bn{সেদ্ধ} \bn{করা} mean? (\textquotedblleft{}To boil\textquotedblright{}.) Say it once more.
\end{tcolorbox}
